% 315_Kfrak_Form_En_ch.tex -- chapter content Doc. 315
% The form of K_frak: additive or multiplicative?
% With Euler's musical spiral (5/7-limit) as control case.

\section{What this is about}

The fractal correction factor appears throughout the corpus in the
additive form
\[
  K_\text{frak}^{\,\text{add}} \;=\; 1 - 100\,\xi \;=\; \frac{74}{75}
  \;=\; 0.98\overline{6},
\]
while a multiplicative--cumulative form would be conceivable:
\[
  K_\text{frak}^{\,\text{mul}} \;=\; (1-\xi)^{100} \;=\; 0.9867543\ldots
\]
The relative distance between the two forms is $8.88\cdot10^{-5}$. The
corpus confirms the \emph{value} $74/75$ several times independently
(A040, A130, A270); whether this also discriminates the \emph{form} --
additive against multiplicative -- had remained unexamined. This
document supplies the examination and answers three questions: Which
corpus location can resolve the two forms? What does recomputation of
the open identities arising there yield? And -- as a structural
question -- \emph{why} should one form be the right one?

Euler's musical spiral serves as the control case. It is the classic
example of a cycle that \emph{never} closes and thereby provides the
counter-probe to the $\xi$ cycle, which closes \emph{exactly} after
$75$ revolutions. The comparison of the two cases supplies, in the end,
the conceptual axis on which the form question hangs: rolled-up versus
unrolled description.

\textbf{Preliminary clarification (avoiding a category error):} The
value $74/75$ is \textbf{[K]} by the A040 derivation; the question of
this document concerns exclusively whether independent witnesses can
\emph{measure} the additive against the multiplicative form -- i.e.\
the status of the form as a measured rather than merely derived
statement.

\section{Control case I: Euler's spiral never closes}

\subsection{The impossibility argument}

\textbf{[Q]} Euler's Tonnetz (\emph{Tentamen novae theoriae musicae},
1739) spans the intervals of just intonation as a lattice of fifths
($3/2$) and major thirds ($5/4$); the octave ($2$) is divided out. The
question \enquote{after how many repetitions does the spiral close?}
has an exact answer: \emph{never}. A closure would require
\[
  2^{x}\cdot 3^{y}\cdot 5^{z} = 1
  \qquad\text{with integers } (x,y,z)\neq(0,0,0),
\]
contradicting the uniqueness of prime factorisation. Equivalently:
$\log 2$, $\log 3$, $\log 5$ are linearly independent over
$\mathbb{Q}$. The orbit lies dense on the octave circle but never meets
its starting point.

\subsection{Near-closures in the 5-limit}

\textbf{[Q]} What exists instead are near-closures whose deficit is in
each case a named comma:

\begin{center}
\begin{tabular}{llll}
\toprule
Path on the Tonnetz & Residual (comma) & Ratio & Cents \\
\midrule
8 fifths $+$ 1 maj.\ third vs.\ 5 octaves & \textbf{schisma} & $32805/32768$ & $\approx 1.95$ \\
4 fifths vs.\ third $+$ 2 octaves & syntonic comma & $81/80$ & $\approx 21.51$ \\
12 fifths vs.\ 7 octaves & Pythagorean comma & $531441/524288$ & $\approx 23.46$ \\
3 maj.\ thirds vs.\ 1 octave & lesser diesis & $128/125$ & $\approx 41.06$ \\
\bottomrule
\end{tabular}
\end{center}

The practically best closure in the 5-limit is the schisma
($\approx 2$ cents, below the hearing threshold); the best-known is
the circle of fifths after $12$ fifths with the Pythagorean comma as
residual.

\subsection{Near-closures in the 7-limit (script result)}

\textbf{[K]} Adding the natural seventh $7/4$ as a third axis (lattice
$\mathbb{Z}^3$ over $3,5,7$), the same impossibility argument applies.
The systematic search (\texttt{euler\_spirale\_7limit.py}, exact
fraction arithmetic, search space $|b|\le 12$ fifths, $|c|\le 6$
thirds, $|d|\le 4$ sevenths; an assertion confirms along the way that
no exact closure ever occurs) yields near-closures an order of
magnitude better than anything in the 5-limit:

\begin{center}
\begin{tabular}{llll}
\toprule
Cents & Comma & Path $3^{b}\,5^{c}\,7^{d}$ & Name \\
\midrule
$0.325$ & $250047/250000$ & $b=-6,\ c=6,\ d=-3$ & landscape comma \\
$0.396$ & $4375/4374$ & $b=-7,\ c=4,\ d=1$ & \textbf{ragisma} \\
$0.721$ & $2401/2400$ & $b=-1,\ c=-2,\ d=4$ & breedsma \\
$1.954$ & $32805/32768$ & $b=-8,\ c=-1,\ d=0$ & schisma (5-limit, for comparison) \\
$5.758$ & $5120/5103$ & $b=-6,\ c=1,\ d=-1$ & hemifamity \\
$7.712$ & $225/224$ & $b=-2,\ c=-2,\ d=1$ & septimal kleisma \\
\bottomrule
\end{tabular}
\end{center}

The ragisma is the simplest sub-cent closure ($7$ fifths down, $4$
thirds and $1$ seventh up miss the starting point by $0.4$ cents); the
coarser 7-limit counterpart of the syntonic comma is Archytas' comma
$64/63$ ($\approx 27$ cents). More prime axes make the lattice denser
-- the spiral approaches itself faster, but never closes in principle.

\subsection{How the spiral is closed after all: temperament}

\textbf{[B]} A closure only arises when commas are \emph{tempered out},
i.e.\ declared equal to $1$. Twelve-tone equal temperament sets the
fifth by decree to exactly $7/12$ of the octave -- it
\emph{rationalises} the step size, and only thereby does the circle of
fifths close. The closed circle is not a property of the pure
intervals but an artefact of rationalisation. This sentence becomes
load-bearing in Section~\ref{sec:eingerollt}.

\section{Control case II: the xi cycle closes exactly}

\textbf{[B]} The $\xi$ cycle stands on the other side of the same
dichotomy. The step size is
\[
  100\,\xi = \frac{400}{30000} = \frac{1}{75}
  \qquad\text{-- exactly rational,}
\]
because $\xi = 4/30000$ is a rational stipulation (A020). A step of
$1/75$ of a cycle closes after exactly $75$ repetitions with zero
residual; $K_\text{frak} = 1 - 1/75 = 74/75$ is the closed path minus
one segment (A040, A050). This is the general torus dichotomy:
rational winding number $\Rightarrow$ closed orbit; irrational
$\Rightarrow$ dense, never-closing orbit.

\begin{center}
\begin{tabular}{lll}
\toprule
 & Musical spiral (5/7-limit) & $\xi$ cycle (FFGFT) \\
\midrule
Step size & $\log_2(3/2),\ \log_2(5/4),\ldots$ irrational & $1/75$ rational \\
Reason & primes mult.\ independent & $\xi$ stipulated rational \\
Behaviour & dense, never meets itself & closes exactly after 75 \\
\enquote{Near-closure} & commas (schisma, ragisma, \ldots) & none needed \\
Verification & numerical only (cents) & exact fraction arithmetic \\
\bottomrule
\end{tabular}
\end{center}

The closure after $75$ is thus no numerical accident but a direct
consequence of the rationality of $\xi$ -- the same rationalisation
that music must impose after the fact through temperament is
stipulated in the framework from the outset.

\section{The form question}
\label{sec:formfrage}

\textbf{[B]} The two forms are identical to first order; their entire
difference is the second binomial term:
\[
  (1-\xi)^{100} = 1 - 100\xi + \binom{100}{2}\xi^2 - \ldots,
  \qquad
  K^{\text{mul}} - K^{\text{add}} = 4950\,\xi^2 - O(\xi^3)
  = 8.76\cdot10^{-5}.
\]
Numerically verified: the binomial term $4950\xi^2 = 8.8\cdot10^{-5}$
covers $99.6\,\%$ of the full distance; the remainder is higher-order.
A corpus location can decide the form only if its own uncertainty is
smaller than this distance -- for power-law locations: smaller than
the correspondingly amplified distance.

\section{Three witnesses}
\label{sec:zeugen}

\subsection{Witness A: the two-route determination (A130)}

\textbf{[K]} A130 determines the scale $E_0$ along two mutually
independent, $\alpha$-free routes: route~1 as the geometric mean
$E_0 = \sqrt{m_e m_\mu}$, route~2 as the construction
$E_0^2 = 4\sqrt2\,m_\mu\,\xi^{-p}$ without $m_e$. Their ratio
\emph{measures} $K_\text{frak}$ without using the factor:
\[
  K_\text{meas} = \bigl(E_0^{(2)}/E_0^{(1)}\bigr)^{-2}.
\]

\textbf{[B]} Two sharpenings relative to A130. First, $m_\mu$ cancels
completely in the ratio; $K_\text{meas}$ depends only on $m_e$
(CODATA, $3\cdot10^{-10}$), $\xi$, the prefactor $4\sqrt2$ and the
exponent $p$. Second, the lever of the exponent is
$\mathrm dK/K = |\ln\xi|\cdot\mathrm dp \approx 8.92\cdot\mathrm dp$:
for witness~A to resolve the form distance, $p$ must be known to
better than $1\cdot10^{-5}$.

\textbf{[K]} Finding:

\begin{center}
\begin{tabular}{lllll}
\toprule
Assumption on $p$ & $K_\text{meas}$ & Res.\ additive & Res.\ multipl. & Verdict \\
\midrule
$p=-0.2679$ (literal, 4 digits) & $0.986220$ & $4.5\cdot10^{-4}$ & $5.4\cdot10^{-4}$ & no form statement \\
$p=-(2-\sqrt3)$ exact & $0.9866533$ & $\mathbf{1.36\cdot10^{-5}}$ & $1.02\cdot10^{-4}$ & \textbf{7.5:1 additive} \\
\bottomrule
\end{tabular}
\end{center}

The four-digit numerical literal $p=-0.2679$ (as it stands in the
check script for A130) does not resolve: its rounding
($\pm5\cdot10^{-5} \to \pm4.5\cdot10^{-4}$ on $K$) exceeds the form
distance fivefold. The identity $p = -(2-\sqrt3)$ would resolve
sharply -- but it is so far neither declared nor derived in the
corpus. Its recomputation follows in Section~\ref{sec:pident}.

\textbf{[K] Register status (R72, incorporated 3~August 2026):} The
register now books the two $E_0$ values no longer as two independent
determinations but as \emph{the same value bare and fractally
corrected} ($7.348$ bare, $7.398$ $\alpha$-anchored; ratio
$1/\sqrt{K_\text{frak}}$ per the A040 rule). From the ratio follows a
determination of the \emph{exponent}: $n = 101.3$, alongside
$n = 100\pm2$ (R67) and $n = 100.27$ (sector ladder, A270). In this
$n$ language the form question reads: additive $K = 1-n\xi$ against
multiplicative $(1-\xi)^{n}$ differ by
$\Delta n_\text{eff} = n(n-1)\xi/2 \approx 0.66$ -- none of the three
$n$ determinations currently reaches this sharpness. The R72 reading
thus further weakens the \enquote{measurement} framing of witness~A:
what remains is the exact $q^{*}$ determination of
Section~\ref{sec:pident}, whose discrimination hangs on the $p$
identity.

\subsection{Witness B: the high-power location (A270)}

\textbf{[K]} A270 carries the identity
$K_\text{frak}^{-36} \approx 16/\pi^2$ with documented agreement
$0.010\,\%$. The exponent $36$ amplifies the form distance to
$36\cdot8.9\cdot10^{-5} = 0.32\,\%$ -- far above this tolerance. The
location is therefore resolution-capable:

\begin{center}
\begin{tabular}{lll}
\toprule
Expression & Value & rel.\ dev.\ from $16/\pi^2 = 1.6211389$ \\
\midrule
$(74/75)^{-36}$ (additive) & $1.6213007$ & $\mathbf{+1.0\cdot10^{-4}}$ \\
$\bigl((1-\xi)^{100}\bigr)^{-36}$ (multipl.) & $1.6161261$ & $-3.1\cdot10^{-3}$ \\
\bottomrule
\end{tabular}
\end{center}

\textbf{Verdict: 31:1 in favour of additive} -- conditional on the
assumption that $16/\pi^2$ is the correct reference. In A270 this is
carried as a coincidence without forward derivation (P35), but
Monte-Carlo secured.

\textbf{[K] Upgrade of the reference through Doc.~314:} The constant is
identifiable corpus-natively. $\pi^2/16$ is the \emph{packing density
of the D4 lattice} (Doc.~314: \enquote{kissing number $24$ against
$8$, packing density $\pi^2/16$ against $\pi^2/32$}) -- hence
\[
  \frac{16}{\pi^2} \;=\; \frac{1}{\Delta(\mathrm{D4})}.
\]
The A270 identity thus reads
$K_\text{frak}^{-36} = 1/\Delta(\mathrm{D4})$: the reciprocal packing
density of precisely the lattice whose structure Doc.~314 verifies in
Hilbert space. The reference is therefore no longer a free-floating
pretty number but a lattice invariant of the framework; P35 narrows
from \enquote{where does the constant come from?} to \enquote{why does
the exponent $36$ couple to the packing density}. The condition of
witness~B becomes considerably weaker thereby, but does not vanish.

\subsection{Witness C: the power form (A040)}

\textbf{[K]} The power form $(D_f^\text{eff}/3)^{D_f^\text{eff}/2}$
with $D_f^\text{eff} = 2.973$ yields $0.986651$: residual additive
$1.6\cdot10^{-5}$, multiplicative $1.0\cdot10^{-4}$ -- tends additive.
However, the four-digit rounding of $D_f^\text{eff}$
($\pm2.5\cdot10^{-4}$ on $K$) exceeds the form distance: not an
independent witness, but consistent with A.

\section{Recomputation of the p identity}
\label{sec:pident}

\textbf{[K]} Since $m_\mu$ cancels, the exponent $q=-p$ can be
computed back \emph{exactly} from the assumed $K$ form (error floor
from $m_e$: $\mathrm dq \approx 3\cdot10^{-11}$):
\[
  q^{*}(K) = \frac{\ln\!\bigl(m_e/(K\cdot4\sqrt2)\bigr)}{\ln\xi},
  \qquad
  \begin{aligned}
    q^{*}_\text{add} &= 0.267950715 && (K=74/75),\\
    q^{*}_\text{mul} &= 0.267960667 && (K=(1-\xi)^{100}),\\
    2-\sqrt3 &= 0.267949192 && (= \tan15^{\circ} = 1/(2+\sqrt3)).
  \end{aligned}
\]

\textbf{[K]} \emph{Asymmetry.} $2-\sqrt3$ lies $7.5\times$ closer to
$q^{*}_\text{add}$. A fair scan of $175$ distinct closed expressions
-- quadratic irrationals $(a+b\sqrt d)/c$ with $|a|,|b|\le9$,
$c\le12$, $d\in\{2,3,5,6,7,10\}$, plus $\pi$, $e$ and $\ln$ forms --
finds in the $\pm3\cdot10^{-6}$ window around $q^{*}_\text{add}$
\emph{exclusively} $2-\sqrt3$, and around $q^{*}_\text{mul}$
\emph{not a single candidate}. If the true exponent has a simple
closed form at all, only the additive reading is compatible with it.

\textbf{[K]} \emph{Significance.} The probability that a random point
lies $\le1.5\cdot10^{-6}$ from \emph{any} candidate of the field is
$\approx1.3\,\%$; from one as low-complexity as $2-\sqrt3$ (the only
one in the window) $\approx10^{-4}$. The proximity thus carries the
same statistical weight as the Monte-Carlo securing of $16/\pi^2$ in
A270. It is not a proof; the $e^2$ warning from A130 (proximity alone
is no argument) applies unchanged.

\textbf{[K]} \emph{Real residual.} Even $q^{*}_\text{add}$ does not
hit $2-\sqrt3$ exactly: residual $+1.52\cdot10^{-6}$ in $q$,
corresponding to $1.36\cdot10^{-5}$ on the $K$ and $m_e$ level
($\approx 7$~eV in $m_e$) -- a factor $\sim45{,}000$ above the
measurement error floor, hence real. The continued fraction
$q^{*}_\text{add} = [0;3,1,2,1,2,1,2,1,2,12,369]$ follows the
periodic pattern of $2-\sqrt3 = [0;3,\overline{1,2}]$ for nine terms
before breaking out; this is an automatic consequence of the
proximity, not additional evidence, but it confirms:
$q^{*}_\text{add}$ is numerically a minutely perturbed quadratic
irrational, not a simple number of its own.

\textbf{[K]} \emph{No correction term found.} The systematic search
for $q = (2-\sqrt3) + c\cdot F(\xi)$ over ten theory families
($\xi$, $\xi^2$, $(100\xi)^2$, $(100\xi)^3$, $\xi\ln(1/\xi)$,
$\xi^{3/2}$, $\xi/\ln(1/\xi)$, $\ln(75/74)$, $\ln^2(75/74)$,
$(2-\sqrt3)\xi$) yields for no $F$ a simple coefficient $c$ (small
fraction or named constant); nor does an alternative exact exponent
exist in the candidate field. The residual remains an open item -- it
is precisely the location A130 addresses as \enquote{a test of the
pole mass against the geometry}; as an order-of-magnitude hint:
$7$~eV lies in the range of QED self-energy differences between mass
definitions.

\textbf{[B]} \emph{Interchangeability.} The $q$ residual and the
prefactor deviation are one degree of freedom in two notations; the
construction fixes only the product $4\sqrt2\cdot\xi^{q}$. Forcing
$q=2-\sqrt3$ exactly, the prefactor would have to be $5.6567774$
instead of $4\sqrt2 = 5.6568542$ ($-1.36\cdot10^{-5}$); no closed
form was found for this either.

\section{Rolled up and unrolled: why the additive form has a reason}
\label{sec:eingerollt}

Sections \ref{sec:zeugen}--\ref{sec:pident} treat the form question as
a measurement problem. This section treats it as a \emph{structural
problem}: the two forms are not two arbitrary computational ans\"atze
but the exact bookkeepings of \emph{two different representations of
the same geometry}. The form question thereby becomes a precise
physical question: In which domain is the correction exactly defined
-- rolled up or unrolled?

\subsection{The two domains}

\textbf{[B]} \emph{Rolled up (T$^4$, winding picture).} On the compact
cycle the natural quantity is the phase or arc length, and phases
\emph{add}. A correction that shortens each winding by the segment
$\varepsilon$ sums linearly over $N$ windings. The closure condition
from A040,
\[
  75\,(1-\varepsilon) = 74 \quad\Rightarrow\quad \varepsilon = \tfrac1{75}
  = 100\xi,
\]
is \emph{exact} in this domain -- as a bookkeeping identity of the
winding, not as an approximation: $75$ shortened revolutions close
precisely when one full revolution is missing in total. The additive
form $K = 1-100\xi$ is therefore not a linearisation of anything; it
is the native, closed bookkeeping of the rolled-up geometry --
representable in exact fraction arithmetic because $\xi$ is rational.

\textbf{[B]} \emph{Unrolled (scale flow, measurement domain).}
Unrolling the cycle in the scale direction turns the winding count
into a run over orders of magnitude (the RG flow from A040,
$\sim17$ decades). In this representation corrections act stepwise on
amplitudes, and stepwise amplitude corrections compose
multiplicatively: $(1-\xi)^{100}\approx e^{-100\xi}$. Were the
unrolled description fundamental, \emph{this} form would be exact and
the additive only its first order.

\subsection{The form distance is the translation error}

\textbf{[B]} By Section~\ref{sec:formfrage} the entire form distance
is the second-order binomial term, $4950\xi^2$. It is thus nothing
other than the error of the translation between the domains at second
order. The form question reads sharply: Which domain carries the
exact definition, which the approximation? Precisely this -- and only
this -- is what witnesses A and B measure.

\subsection{The Euler parallel makes the assignment unambiguous}

\textbf{[B]} The musical spiral provides the control case in which the
domain assignment is \emph{uncontested} -- because there,
measurement takes place. What ear and instrument reach are
frequencies, i.e.\ quantities of the unrolled scale domain; intervals
compose multiplicatively there ($3/2\cdot3/2\cdot\ldots$), and only
the logarithm rolls them onto the octave circle -- whereby the
irrationality enters and prevents closure in principle. The $\xi$
cycle, conversely, is \emph{defined} in the rolled-up domain: from
the outset a circle structure with rational winding number, without
two multiplicatively defined objects that would first have to be made
comparable logarithmically.

\begin{center}
\begin{tabular}{lll}
\toprule
 & Music (Euler) & FFGFT ($\xi$ cycle) \\
\midrule
Measured quantity & frequency (unrolled) & --- (defined rolled up) \\
Composition & intervals multiply & winding phases add \\
Bridge to the circle & logarithm (forced) & none needed \\
Step size & $\log_2(3/2)$ irrational & $1/75$ rational \\
Closure & never (commas only) & exact after 75 \\
\bottomrule
\end{tabular}
\end{center}

The assignment is complete: multiplicative composition belongs to the
unrolled measurement/scale domain, additive composition to the
rolled-up winding domain.

\subsection{The corpus already contains the fork: case B against case C}

\textbf{[K]} The assignment does not stand beside the corpus but
within it. The closure fork (Docs.~295/313, cited in Doc.~314) carries
three cases: case~A ($d=0$, closure after one revolution), case~B
($\xi$ frozen: rational rotation number $1-d = 74/75$; because
$\gcd(74,75)=1$, closure after exactly $75$ revolutions) and case~C
($\xi$ runs via the recursion: \enquote{the cumulated defect becomes
logarithmic and the winding an equiangular spiral} -- no closure).
This is precisely the form question as physics: case~B \emph{is} the
additive bookkeeping (linear defect, rational winding, exact
closure), case~C \emph{is} the multiplicative--logarithmic
composition -- and the equiangular spiral is the same mathematical
picture as the Euler spiral of control case~I. The two forms of
$K_\text{frak}$ are thus the bookkeepings of the two branches of the
fork; the question \enquote{additive or multiplicative?} is the
question of which branch carries the correction.

Doc.~314 (Ch.~D2) additionally sharpens the domain boundary in the
terminology used here: \emph{in the rolled-up state there are no
fractal defects; the defect $d = 100\xi$ accumulates only unrolled},
and $K_\text{frak}$ is \emph{a property of the traversal}, entering
via the bridge constants ($E_0$, $v$), not via the local operator.
The additive form is accordingly the bookkeeping of the \emph{closed}
traversal with frozen $\xi$ (case~B); the multiplicative would be the
bookkeeping of the \emph{running} accumulation (case~C). All corpus
uses of $K_\text{frak}$ (A040 rule, R70 assignment, R72) book
case~B.

\subsection{What the witnesses then mean}

\textbf{[B]} In this reading, witnesses A and B do not measure a
computational detail but the domain membership of the correction --
and they answer in agreement: the exact bookkeeping is the
\emph{rolled-up} one. The multiplicative form is the approximation a
naive observer of the unrolled scale narrative would write down --
correct to first order, wrong from $4950\xi^2$ on.

This also resolves the apparent tension within A040: the document
contains both -- the closure comma (rolled-up, exact reading) and the
cumulative RG flow (unrolled derivation of the factor $100$). Both
are consistent because they agree to first order; by the witness
evidence, the rolled-up one is exact. The RG flow derives \emph{why}
the correction has magnitude $100\xi$; the winding geometry fixes
\emph{how} it composes. The naive RG argument for the multiplicative
form is thereby placed rather than refuted: it is the correct
argument of the unrolled domain -- and that domain, by the witness
evidence, is not the defining one.

\subsection{Connection to the foundational architecture}

\textbf{[B]} The assignment \enquote{rolled up $=$ exact} is not an
ad-hoc rescue of the additive form but congruent with the
foundational architecture of the framework: the compact geometry
(T$^4$, rational $\xi$ stipulation) is primary; scale-flow narratives
are derived descriptions. That the exact arithmetic of the corpus
works throughout with fractions is the same statement in
methodological form: exact fractions exist only in the rolled-up
domain -- the unrolled one knows only approximations with $\xi^2$
residuals. In this respect a consistently rolled-up framework
\emph{must} carry the additive form; were the measurement to come out
multiplicative, not only $K_\text{frak}$ would be affected but the
primacy claim of the geometry itself.

\section{Status and falsification line}

\begin{keyresult}[Bookkeeping status]
$K_\text{frak} = 1-100\xi = 74/75$: value \textbf{[K]} (A040). Form
additive: \textbf{[B]}, twice conditionally confirmed --
condition~1 ($p$ identity $2-\sqrt3$) unchanged open; condition~2
($16/\pi^2$ reference) substantially weakened by Doc.~314, since
$16/\pi^2 = 1/\Delta(\mathrm{D4})$ is identified corpus-natively and
only the coupling to the exponent~$36$ remains open (P35 narrowed).
Structurally the additive form is the bookkeeping of fork branch
case~B (frozen $\xi$, rational winding), which all corpus uses book;
the multiplicative alternative $(1-\xi)^{100}$ is the bookkeeping of
the spiral branch case~C, is consistently disfavoured by both
conditional witnesses and supported by none.
\end{keyresult}

\textbf{[B]} \emph{Why nevertheless not [K]-measured:} Each of the two
sharp discriminations hangs on a still-open premise -- witness~A on
the derivation of the identity $p=-(2-\sqrt3)$ (with real residual
$1.5\cdot10^{-6}$), witness~B on the status of $16/\pi^2$ as
reference (P35, narrowed). If \emph{one} of the two premises later
turns out correct, the form statement becomes unconditional.
Conversely, the multiplicative form remains logically possible as
long as both premises are open: then the $16/\pi^2$ proximity and the
$2-\sqrt3$ proximity would be two independent coincidences -- and the
naive RG argument of the unrolled domain would have been right.

\textbf{[B]} \emph{Falsification line.} The question would be
decidable unconditionally at a location where $K$ stands at high
power against a \emph{measurement} rather than a coincidence
constant. The candidate already stands in the corpus: the baryon
prediction from A270 (sector exponent $+2$, i.e.\ $K^{38}$ level,
form distance there $\approx0.34\,\%$). A baryon mass check of better
accuracy would decide additive against multiplicative without
conditions -- and at the same time test the domain assignment of
Section~\ref{sec:eingerollt}: were it to come out multiplicative, not
only the form choice would be corrected; the correction would live in
the scale domain, against the winding reading of the closure comma.
The structural argument thus raises the stakes; it binds the form
question to the primacy question.

\section{Open items}

\begin{itemize}
  \item \textbf{P-315-1:} Derivation or declaration of the identity
    $p=-(2-\sqrt3)$ in A130/route~2 (currently only a four-digit
    numerical literal in the check script). $2-\sqrt3=\tan15^{\circ}$
    has plausible geometric flavour (one-twelfth angle) but no corpus
    derivation. Per R72 this is the only way to make witness~A
    resolution-capable.
  \item \textbf{P-315-2:} The real residual $1.36\cdot10^{-5}$ on the
    $K$/$m_e$ level ($\approx7$~eV): a correction term of the theory,
    a pole-mass effect, or a hint at a different exact closure form?
    None of the ten tested $F(\xi)$ families yields a simple
    coefficient. Consistent with R72: the $n$ determination from the
    $E_0$ ratio ($101.3$) lies above $100$; the excess is the same
    item.
  \item \textbf{P35 (narrowed by Doc.~314):} With
    $16/\pi^2 = 1/\Delta(\mathrm{D4})$ the constant is identified as
    the reciprocal D4 packing density; what remains open is the
    forward derivation of the coupling
    $K_\text{frak}^{-36} = 1/\Delta(\mathrm{D4})$, in particular the
    role of the bulk exponent~$36$.
  \item \textbf{P-315-3 ($n$ sharpness):} In $n$ language, form
    discrimination requires $\Delta n \le 0.66$ ($= n(n-1)\xi/2$).
    The three $n$ routes from R72 ($101.3$; $100\pm2$; $100.27$) do
    not reach this; the sector ladder ($100.27$) is the most
    promising candidate for sharpening.
  \item \textbf{Sharpening A040:} state $D_f^\text{eff}$ with more
    digits (needed: better than $\pm1\cdot10^{-4}$) so that
    witness~C becomes resolution-capable.
  \item \textbf{Baryon location (A270):} mass check at the $K^{38}$
    level as unconditional form test -- simultaneously a test of the
    fork assignment case~B/C.
\end{itemize}

\section{Verification}

All numbers of this document are reproducible by four scripts (pure
standard library; exact fraction arithmetic where possible):

\begin{itemize}
  \item \texttt{python/315\_Skripte/euler\_spirale\_7limit.py} --
    comma search in the 5/7-limit, exact; assertion against exact
    closure.
  \item \texttt{python/315\_Skripte/pruefrechnung\_kfrak\_form.py} --
    witnesses A/B/C with verdict logic (uncertainty against the
    effective, possibly amplified form distance).
  \item \texttt{python/315\_Skripte/pruefrechnung\_p\_identitaet.py} --
    exact $q^{*}$ determination, fairness scan, residual and
    prefactor cross-check.
  \item \texttt{python/315\_Skripte/pruefrechnung\_rest\_0p1xi.py} --
    correction-term families, candidate statistics (random
    expectation), continued-fraction diagnosis.
\end{itemize}

Inputs: $\xi=4/30000$ exact; $m_e = 0.51099895069$~MeV,
$m_\mu = 105.6583755$~MeV (CODATA~2022).
